% Insert in the supplementary Discussion. Replace its existing P2 bootstrap
% sentence with the first paragraph; keep the root's F4/G2 and matcher discussion.

P2 contains one retained run per task and configuration. Its task-cluster bootstrap intervals describe variation across the sampled task mix, not stochastic variation from independent runs of the same task. The separate P1b study supplies three repeats on each of four tasks under each configuration; these 48 runs do not provide repeated measurements for every task in the 50-task P2 study. The complete organized P2 table covers all 200 runs and their three score components. Some Hard-task accuracy targets concern qualitative channel identity, so the reported gain includes pathway identification and evidence completion as well as numerical barriers. Repeating the full task set would characterize run-to-run variability and the frequency with which each research route is selected.

Cloudra's default preorganization samples 32 rigid relative placements, applies local MMFF94 relaxation, and selects candidates by energy and clustering. This is a tractable candidate-generation procedure, not an exhaustive search for the global minimum of every reactant arrangement or intermediate. CREST is available through a separate optional backend; the 200-run study does not establish uniform CREST coverage. Highly flexible systems may require broader sampling and more quantum-chemical calculations. Further work should measure conformer coverage and improve the allocation of calculations among competing candidates.

The benchmark compares electronic energy differences. It does not add zero-point, thermal/entropic, or standard-state corrections, and it does not quantify the free-energy contribution of solvation. This scope matters for NI5, specified in weakly acidic aqueous solution, and NU5, specified in dichloromethane at high concentration. Their scores assess the prescribed electronic-energy references rather than experimental solution-phase free-energy barriers. No solvent-sensitivity comparison was performed, so the present results do not determine the magnitude or direction of a solvent effect.

The atom-mapping coefficients and ambiguity window, and the decay, depth, and congestion settings in the branch-crowding descriptor, were fixed across the workflow comparisons. We did not vary them in a parameter sweep or measure each setting's separate contribution. The imaginary-mode alignment ratio of 0.5 is advisory; the stationary-point and IRC/endpoint evidence determine scientific acceptance. All P2 configurations used the archived API identifier \texttt{deepseek-v4-flash} with the recorded request settings. The archive identifies that API configuration rather than an immutable provider-side weight/build snapshot. Comparisons with fixed model snapshots and additional backbones would test how the workflow gains carry across model versions.
